\chapter{Conclusion et Perspectives}

\section{Conclusion}

La conception de cette plateforme web souveraine répond à un besoin stratégique double et majeur pour l'ANTIC : d'une part, l'automatisation du \textbf{filtrage dynamique des données Excel et la création de feuilles enfants} directement au sein des classeurs sources, et d'autre part, la \textbf{conversion sécurisée et la mise en forme automatisée vers des rapports Word normalisés}. En substituant un processus manuel, fragmenté et risqué par une application centralisée, performante et souveraine, l'institution franchit un cap décisif dans la modernisation de son ingénierie documentaire.

L'architecture N-Tiers retenue, associant la simplicité de Python FastAPI, openpyxl et python-docx à la légèreté de HTML5/JS/Bootstrap et SQLite/PostgreSQL, offre toutes les garanties de sécurité, d'auditabilité et de scalabilité requises par les standards de l'État.

\section{Perspectives d'évolution}

La modularité de la solution proposée permet d'envisager plusieurs extensions futures :
\begin{itemize}
    \item \textbf{Génération directe au format PDF :} Ajout d'un module d'export immédiat au format PDF signé électroniquement par le certificat de l'ANTIC ;
    \item \textbf{Connecteurs automatisés :} Connexion directe avec les systèmes d'information internes pour l'extraction automatique des données sans soumission manuelle du fichier Excel ;
    \item \textbf{Module d'IA décisionnelle :} Analyse sémantique automatique des données converties pour générer des résumés exécutifs intelligents en tête de rapport.
\end{itemize}
